% !TEX program = xelatex
% Самодостаточный конспект. Компилировать XeLaTeX/LuaLaTeX (tectonic 07-znakoperemennye-ryady.tex).
\documentclass[11pt,a4paper]{article}
\newcommand{\coursename}{Математический анализ}
\newcommand{\rightheader}{§7. Знакопеременные ряды}
\newcommand{\doctitle}{Знакопеременные ряды}
% ==== Общая преамбула конспектов (собирается в каждый .tex как самодостаточная) ====
% Перед этим файлом должны быть определены: \documentclass и макросы
%   \coursename  — название курса (левый колонтитул, подзаголовок),
%   \rightheader — правый колонтитул (напр. «§2. Рациональные дроби»),
%   \doctitle    — заголовок раздела на титуле.
\usepackage{fontspec}
\setmainfont{cmunrm.otf}[
  BoldFont       = cmunbx.otf,
  ItalicFont     = cmunti.otf,
  BoldItalicFont = cmunbi.otf]
\usepackage{polyglossia}
\setdefaultlanguage{russian}
\setotherlanguage{english}

\usepackage{amsmath,amssymb,amsthm,mathtools}
\usepackage[a4paper,margin=2.2cm]{geometry}
\usepackage{enumitem}
\usepackage{graphicx}
\usepackage{array}
\usepackage{xcolor}
\definecolor{accent}{HTML}{1F6FEB}
\definecolor{rulecol}{HTML}{D0D7DE}
\usepackage[colorlinks=true,linkcolor=accent,urlcolor=accent,citecolor=accent]{hyperref}
\usepackage{titlesec}
\usepackage{fancyhdr}

% ----- Колонтитулы -----
\pagestyle{fancy}
\fancyhf{}
\fancyhead[L]{\small\itshape \coursename}
\fancyhead[R]{\small\itshape \rightheader}
\fancyfoot[C]{\thepage}
\renewcommand{\headrulewidth}{0.4pt}
\renewcommand{\headrule}{\hbox to\headwidth{\color{rulecol}\leaders\hrule height \headrulewidth\hfill}}

% ----- Заголовки -----
\titleformat{\section}{\large\bfseries\color{accent}}{\thesection.}{0.6em}{}
\titleformat{\subsection}{\bfseries}{\thesubsection.}{0.5em}{}

% ----- Теоремные окружения (стиль конспекта) -----
\theoremstyle{definition}
\newtheorem{definition}{Определение}[section]
\newtheorem{example}[definition]{Пример}
\newtheorem{remark}[definition]{Замечание}
\theoremstyle{plain}
\newtheorem{theorem}[definition]{Теорема}
\newtheorem{lemma}[definition]{Лемма}
\newtheorem{corollary}[definition]{Следствие}
\newtheorem{property}[definition]{Свойство}
\newtheorem{criterion}[definition]{Критерий}
\renewcommand{\proofname}{Доказательство}

% ----- Операторы и сокращения -----
\DeclareMathOperator{\tg}{tg}
\DeclareMathOperator{\ctg}{ctg}
\DeclareMathOperator{\arctg}{arctg}
\DeclareMathOperator{\arcctg}{arcctg}
\DeclareMathOperator{\sh}{sh}
\DeclareMathOperator{\ch}{ch}
\DeclareMathOperator{\Th}{th}
\DeclareMathOperator{\cth}{cth}
\DeclareMathOperator{\Const}{const}
\DeclareMathOperator{\sign}{sign}
\DeclareMathOperator{\rang}{rang}
\DeclareMathOperator{\diam}{diam}
\DeclareMathOperator{\Span}{span}
\DeclareMathOperator{\Ker}{Ker}
\DeclareMathOperator{\Img}{Im}
\DeclareMathOperator{\tr}{tr}
\newcommand{\R}{\mathbb{R}}
\newcommand{\C}{\mathbb{C}}
\newcommand{\N}{\mathbb{N}}
\newcommand{\Z}{\mathbb{Z}}
\newcommand{\Q}{\mathbb{Q}}
\newcommand{\dx}{\,dx}
\newcommand{\dt}{\,dt}
\newcommand{\du}{\,du}
\newcommand{\eps}{\varepsilon}
\renewcommand{\Re}{\operatorname{Re}}
\renewcommand{\Im}{\operatorname{Im}}
\renewcommand{\le}{\leqslant}
\renewcommand{\ge}{\geqslant}
\newcommand{\scal}[2]{\left( #1,\,#2\right)}
\DeclarePairedDelimiter{\abs}{\lvert}{\rvert}
\DeclarePairedDelimiter{\norm}{\lVert}{\rVert}

\begin{document}
\begin{center}
  {\LARGE\bfseries \doctitle}\\[4pt]
  {\large\itshape \coursename}
\end{center}
\vspace{6pt}

\section{Абсолютная и условная сходимость}

Пусть теперь члены ряда $\displaystyle\sum_{k=1}^\infty a_k$ могут иметь любой знак.

\begin{definition}
Ряд $\sum a_k$ \emph{сходится абсолютно}, если сходится ряд из модулей $\sum\abs{a_k}$.
\end{definition}

\begin{theorem}[о связи сходимостей]
Абсолютно сходящийся ряд сходится.
\end{theorem}

\begin{proof}
Пусть $\sum\abs{a_k}$ сходится. По критерию Коши $\forall\eps>0\ \exists N\ \forall n>N\ \forall
p$: $\abs{a_{n+1}}+\dots+\abs{a_{n+p}}<\eps$. Тогда по неравенству треугольника
$\abs{a_{n+1}+\dots+a_{n+p}}\le\abs{a_{n+1}}+\dots+\abs{a_{n+p}}<\eps$, т.е. условие Коши
выполнено и для $\sum a_k$.
\end{proof}

\begin{definition}
Ряд \emph{сходится условно}, если он сходится, но не абсолютно.
\end{definition}

\section{Знакочередующиеся ряды. Признак Лейбница}

\begin{definition}
Ряд вида $\displaystyle\sum_{k=1}^\infty(-1)^{k+1}a_k=a_1-a_2+a_3-a_4+\dots$ с $a_k>0$ называют
\emph{знакочередующимся}.
\end{definition}

\begin{theorem}[Лейбниц]
Если $a_k>0$ монотонно убывают и $a_k\to0$, то знакочередующийся ряд
$\sum(-1)^{k+1}a_k$ сходится. При этом его сумма $0<S<a_1$, а остаток
$\displaystyle R_n=\sum_{k=n+1}^\infty(-1)^{k+1}a_k$ оценивается как $\abs{R_n}\le a_{n+1}$.
\end{theorem}

\begin{proof}
Чётные частичные суммы $S_{2n}=(a_1-a_2)+(a_3-a_4)+\dots+(a_{2n-1}-a_{2n})$ не убывают (каждая
скобка $\ge0$). С другой стороны, $S_{2n}=a_1-(a_2-a_3)-\dots-(a_{2n-2}-a_{2n-1})-a_{2n}<a_1$,
т.е. $\{S_{2n}\}$ ограничена сверху; по теореме Вейерштрасса $S_{2n}\to S$. Наконец
$S_{2n+1}=S_{2n}+a_{2n+1}\to S+0=S$, поэтому и $S_n\to S$. Оценка остатка следует из того, что
$R_n$ — снова лейбницевский ряд с первым членом $a_{n+1}$.
\end{proof}

\begin{example}
$\displaystyle\sum_{k=1}^\infty\frac{(-1)^{k+1}}{k}=1-\frac12+\frac13-\frac14+\dots=\ln2$ сходится
по признаку Лейбница, но \emph{условно} (ряд из модулей — гармонический, расходится).
\end{example}

\section{Перестановки членов}

\begin{theorem}[Риман]
Если ряд сходится условно, то для любого $\ell\in\R$ (в том числе $\ell=\pm\infty$) можно так
переставить его члены, что новый ряд будет сходиться к $\ell$.
\end{theorem}

\begin{theorem}[Дирихле]
Если ряд сходится абсолютно, то любая его перестановка сходится к той же сумме.
\end{theorem}

\begin{remark}
Иллюстрация теоремы Римана: переставив члены ряда
$\sum\frac{(-1)^{k+1}}{k}=S$ как
$\bigl(1-\tfrac12-\tfrac14\bigr)+\bigl(\tfrac13-\tfrac16-\tfrac18\bigr)+\dots$, получим ряд с
суммой $\tfrac12 S$ — перестановка изменила сумму, что невозможно для абсолютно сходящихся рядов.
\end{remark}

\end{document}
